% Options for packages loaded elsewhere
\PassOptionsToPackage{unicode}{hyperref}
\PassOptionsToPackage{hyphens}{url}
\PassOptionsToPackage{dvipsnames,svgnames,x11names}{xcolor}
%
\documentclass[
  10pt,
]{article}
\usepackage{amsmath,amssymb}
\usepackage{iftex}
\ifPDFTeX
  \usepackage[T1]{fontenc}
  \usepackage[utf8]{inputenc}
  \usepackage{textcomp} % provide euro and other symbols
\else % if luatex or xetex
  \usepackage{unicode-math} % this also loads fontspec
  \defaultfontfeatures{Scale=MatchLowercase}
  \defaultfontfeatures[\rmfamily]{Ligatures=TeX,Scale=1}
\fi
\usepackage{lmodern}
\ifPDFTeX\else
  % xetex/luatex font selection
  \setmainfont[]{DejaVu Serif}
  \setsansfont[]{DejaVu Sans}
  \setmonofont[]{DejaVu Sans Mono}
  \setmathfont[]{Latin Modern Math}
\fi
% Use upquote if available, for straight quotes in verbatim environments
\IfFileExists{upquote.sty}{\usepackage{upquote}}{}
\IfFileExists{microtype.sty}{% use microtype if available
  \usepackage[]{microtype}
  \UseMicrotypeSet[protrusion]{basicmath} % disable protrusion for tt fonts
}{}
\makeatletter
\@ifundefined{KOMAClassName}{% if non-KOMA class
  \IfFileExists{parskip.sty}{%
    \usepackage{parskip}
  }{% else
    \setlength{\parindent}{0pt}
    \setlength{\parskip}{6pt plus 2pt minus 1pt}}
}{% if KOMA class
  \KOMAoptions{parskip=half}}
\makeatother
\usepackage{xcolor}
\usepackage[margin=2.0cm]{geometry}
\usepackage{longtable,booktabs,array}
\usepackage{calc} % for calculating minipage widths
% Correct order of tables after \paragraph or \subparagraph
\usepackage{etoolbox}
\makeatletter
\patchcmd\longtable{\par}{\if@noskipsec\mbox{}\fi\par}{}{}
\makeatother
% Allow footnotes in longtable head/foot
\IfFileExists{footnotehyper.sty}{\usepackage{footnotehyper}}{\usepackage{footnote}}
\makesavenoteenv{longtable}
\setlength{\emergencystretch}{3em} % prevent overfull lines
\providecommand{\tightlist}{%
  \setlength{\itemsep}{0pt}\setlength{\parskip}{0pt}}
\setcounter{secnumdepth}{-\maxdimen} % remove section numbering
% doc_header.tex -- LaTeX preamble additions used by build_docs.sh for every
% Markdown -> LaTeX -> PDF conversion (supplements, README, submission files).
\usepackage{etoolbox}
% longtable may not be loaded by pandoc when a document has no table
\usepackage{longtable}
% bold vectors: \bm -> \symbf under unicode-math (XeLaTeX), bm package otherwise
\makeatletter
\@ifpackageloaded{unicode-math}{\providecommand{\bm}[1]{\symbf{#1}}}{\usepackage{bm}}
\makeatother
% smaller type in long tables so that wide audit tables fit the text width;
% inside a table a line break is allowed after an underscore of plain text
% (pandoc writes it as \_), so that long identifiers wrap inside their column
% rows are set a little apart, so that symbols with indices do not touch
\AtBeginEnvironment{longtable}{\footnotesize\renewcommand{\arraystretch}{1.45}%
  \let\sptunderscore\_\def\_{\sptunderscore\allowbreak}}
\setlength{\LTcapwidth}{\textwidth}
% allow line breaks after slashes and hyphens inside table cells
\usepackage{xurl}
% code blocks: break over-long lines instead of running into the margin.
% Plain ``` blocks are typeset by pandoc as {verbatim}, highlighted ones
% (```bash) as {Highlighting}; both are re-declared as breakable fancyvrb
% environments.  The continuation line is marked with a small hooked arrow.
\usepackage{fvextra}
\RecustomVerbatimEnvironment{verbatim}{Verbatim}{breaklines,breakanywhere,fontsize=\small}
\makeatletter
\@ifundefined{Highlighting}{}{%
  \RecustomVerbatimEnvironment{Highlighting}{Verbatim}{breaklines,breakanywhere,fontsize=\small,commandchars=\\\{\}}}
\makeatother
% inline code (pandoc writes it as \texttt{...}): never hyphenate -- an inserted
% hyphen could be read as part of a file name -- but allow a line break after
% an underscore, slash, comma, colon or equals sign, so that long paths and
% column lists wrap inside the text width and inside table columns.
% Inside a formula (e.g. the name of a notebook symbol in an equation) the same
% command must first leave math mode, which \text of amsmath does.
\ExplSyntaxOn
\tl_new:N \l__spt_code_tl
\cs_new_protected:Npn \__spt_texttt:n #1
  {
    \tl_set:Nn \l__spt_code_tl {#1}
    \regex_replace_all:nnN { ([\/\,\:\=]) } { \1 \c{allowbreak} } \l__spt_code_tl
    \regex_replace_all:nnN { \c{_} } { \c{_} \c{allowbreak} } \l__spt_code_tl
    \group_begin: \ttfamily \hyphenchar \font = -1 \scan_stop: \l__spt_code_tl \group_end:
  }
\RenewDocumentCommand \texttt { m }
  {
    \mode_if_math:TF { \text { \__spt_texttt:n {#1} } } { \__spt_texttt:n {#1} }
  }
\ExplSyntaxOff
% A figure is written in the Markdown source as an image paragraph followed by
% a caption paragraph ("**Figure S1.** ..."), so that the caption is also
% visible on GitHub.  An image that starts a paragraph is centered.
\usepackage{graphicx}% pandoc loads it only for documents that contain an image
\NewCommandCopy\sptincludegraphics\includegraphics
\RenewDocumentCommand\includegraphics{O{} m}{%
  \ifvmode\begin{center}\sptincludegraphics[#1]{#2}\end{center}%
  \else\sptincludegraphics[#1]{#2}\fi}
% keep figures where they are referenced
\usepackage{float}
\floatplacement{figure}{H}
% running header with the document title
\usepackage{fancyhdr}
\pagestyle{fancy}
\fancyhf{}
\fancyhead[L]{\small PROJECT spacetimes --- documentation and provenance}
\fancyhead[R]{\small\thepage}
\renewcommand{\headrulewidth}{0.4pt}
% 16x16 matrices: tighter columns so that a full matrix fits the text width
\setlength{\arraycolsep}{3.2pt}
% The documents of this project contain long tagged symbols, file paths and
% identifiers that cannot be hyphenated.  A paragraph is therefore allowed to
% have looser lines before TeX gives up and lets a line stick out of the margin.
\AtBeginDocument{\tolerance=2000 \emergencystretch=5em \hbadness=2000}
\ifLuaTeX
  \usepackage{selnolig}  % disable illegal ligatures
\fi
\IfFileExists{bookmark.sty}{\usepackage{bookmark}}{\usepackage{hyperref}}
\IfFileExists{xurl.sty}{\usepackage{xurl}}{} % add URL line breaks if available
\urlstyle{same}
\hypersetup{
  pdftitle={S2 --- Is the hidden-axis singlet superfluid the Higgs field? An assessment},
  pdfauthor={Prepared for Patrick Lee Nash, Ph.D.},
  colorlinks=true,
  linkcolor={blue},
  filecolor={Maroon},
  citecolor={blue},
  urlcolor={blue},
  pdfcreator={LaTeX via pandoc}}

\title{S2 --- Is the hidden-axis singlet superfluid the Higgs field? An
assessment}
\usepackage{etoolbox}
\makeatletter
\providecommand{\subtitle}[1]{% add subtitle to \maketitle
  \apptocmd{\@title}{\par {\large #1 \par}}{}{}
}
\makeatother
\subtitle{Supplement to both manuscripts of PROJECT spacetimes: the
conditions that the hypothesis must meet, which of them the 7+1 and the
4+4 model meet, and what remains open}
\author{Prepared for Patrick Lee Nash, Ph.D.}
\date{November 12, 2026}

\begin{document}
\maketitle

\hypertarget{purpose-and-result}{%
\section{0. Purpose and result}\label{purpose-and-result}}

The task states the hypothesis that the superfluid singlet state of
Supplement S1 --- the condensate of the pair operator along the hidden
axis, \(\Delta^0\propto\langle\psi^TC^{-}\gamma^0\psi\rangle\) ---
\emph{is the Higgs field}, and asks for a discussion with provenance.
This supplement lists what the Higgs field of the Standard Model is
required to do (Sec. 1), tests the condensate of each spacetime against
every requirement (Secs. 2 to 5), and states the result (Sec. 6):

\begin{itemize}
\tightlist
\item
  \textbf{7+1.} The condensate has the right transformation properties:
  it is a scalar in ordinary spacetime, and its symmetry-breaking
  pattern, \(SO(4)\to SO(3)\) for the rotations of the four extra space
  axes, is exactly the pattern of the Higgs vacuum. This is a
  non-trivial necessary condition, and it is met. None of the dynamical
  conditions has been shown, and the simplest estimates of two of them
  are unfavourable.
\item
  \textbf{4+4.} For fields that do not depend on the extra coordinates
  the condensate has the same pattern, but under a \emph{hidden}
  symmetry that is not a transformation of the spacetime and that the
  geometry cannot turn into a gauge symmetry. Under the symmetry that
  the geometry of 4+4 does provide --- the rotations of the three
  additional times --- the hidden-axis condensate is a singlet and
  cannot break the electroweak symmetry. The hypothesis has no geometric
  basis in 4+4.
\end{itemize}

The symbols are those of ``dirac matrices'' (cited as DM) and of S1. The
calculations referred to as ``S4'' are those of Supplement S4 (the
effective action of Cornwall, Jackiw and Tomboulis).

\begin{center}\rule{0.5\linewidth}{0.5pt}\end{center}

\hypertarget{what-the-higgs-field-does-and-the-resulting-conditions}{%
\section{1. What the Higgs field does, and the resulting
conditions}\label{what-the-higgs-field-does-and-the-resulting-conditions}}

In the Standard Model the Higgs field \(H\) is a doublet of complex
scalar fields under the weak-isospin group \(SU(2)_L\), with weak
hypercharge \(Y=\tfrac12\). Its potential is minimal at
\(|H|=v_{\rm EW}/\sqrt2\), \(v_{\rm EW}=(\sqrt2G_F)^{-1/2}=246.22\) GeV.
This vacuum value breaks \(SU(2)_L\times U(1)_Y\) to the \(U(1)\) of
electromagnetism, gives the \(W\) and \(Z\) bosons their masses with the
ratio \(\rho\equiv M_W^2/(M_Z^2\cos^2\theta_W)=1\), and gives the quarks
and charged leptons their masses through Yukawa couplings. The one
remaining scalar particle, \(h\), has the mass 125.20 GeV (Particle Data
Group) and couplings proportional to the masses of the particles it
couples to, measured to within about ten per cent.

Two structural facts are essential for what follows.

\emph{The custodial symmetry.} The four real components of \(H\) form a
vector of \(SO(4)\), which is locally \(SU(2)_L\times SU(2)_R\); the
Higgs potential depends only on its length. The vacuum value points in
one direction and breaks \(SO(4)\) to \(SO(3)\). That residual \(SO(3)\)
is what protects \(\rho=1\) (Sikivie, Susskind, Voloshin and Zakharov).
In the language of representations, \(H\) is a \emph{real} \((2,2)\).

\emph{Composite Higgs fields.} The idea that the Higgs field is not
elementary but a condensate of fermion pairs is old: it is the idea of
Nambu and Jona-Lasinio applied to the electroweak scale (Bardeen, Hill
and Lindner, with top quarks). The present hypothesis is of this kind,
with one difference: the pair is made of two fermions (it carries
fermion number 2), not of a fermion and an antifermion.

A condensate that is to \emph{be} the Higgs field must satisfy:

\begin{longtable}[]{@{}
  >{\raggedright\arraybackslash}p{(\columnwidth - 4\tabcolsep) * \real{0.0385}}
  >{\raggedright\arraybackslash}p{(\columnwidth - 4\tabcolsep) * \real{0.3590}}
  >{\raggedright\arraybackslash}p{(\columnwidth - 4\tabcolsep) * \real{0.6026}}@{}}
\toprule\noalign{}
\begin{minipage}[b]{\linewidth}\raggedright
\end{minipage} & \begin{minipage}[b]{\linewidth}\raggedright
Condition
\end{minipage} & \begin{minipage}[b]{\linewidth}\raggedright
Why
\end{minipage} \\
\midrule\noalign{}
\endhead
\bottomrule\noalign{}
\endlastfoot
C1 & scalar under the Lorentz group of ordinary spacetime & the vacuum
is Lorentz invariant \\
C2 & non-zero in empty space & the vacuum contains no medium \\
C3 & a doublet of \(SU(2)_L\) with a neutral component & to break the
electroweak symmetry and leave electromagnetism \\
C4 & symmetry breaking \(SO(4)\to SO(3)\) & \(\rho=1\) \\
C5 & scale 246 GeV; a scalar particle of 125 GeV with the observed
couplings & measured \\
C6 & couplings that give the quarks and leptons their masses &
measured \\
C7 & no further light or charged particles that would have been seen &
measured \\
\end{longtable}

\begin{center}\rule{0.5\linewidth}{0.5pt}\end{center}

\hypertarget{condition-c1-a-scalar-in-ordinary-spacetime-met-in-both-spacetimes}{%
\section{2. Condition C1: a scalar in ordinary spacetime --- met in both
spacetimes}\label{condition-c1-a-scalar-in-ordinary-spacetime-met-in-both-spacetimes}}

The pair operator \(\Delta^0=\psi^TC^{-}\gamma^0\psi\) is the component
along \(x0\) of a vector of the eight-dimensional spacetime. It is
invariant under all Lorentz transformations that do not involve \(x0\)
(DM, Sec. 3.3; checks \texttt{44-K3}, \texttt{71-K3}): under the Lorentz
group of the seven coordinates \(x1,\dots,x7\), which contains the
Lorentz group of ordinary spacetime \((x1,x2,x3;x4)\). This is the
property for which the hidden axis is needed: a scalar pair operator of
the full eight-dimensional group does not exist (DM, Theorem 4), but a
vector along a dimension that ordinary physics does not see is a scalar
for ordinary physics.

\begin{center}\rule{0.5\linewidth}{0.5pt}\end{center}

\hypertarget{condition-c2-a-condensate-in-empty-space-not-established}{%
\section{3. Condition C2: a condensate in empty space --- not
established}\label{condition-c2-a-condensate-in-empty-space-not-established}}

The superfluid of S1 forms in a \emph{medium} of fermions: the gap
equation (S1, Eq. (5.2)) has a solution for an arbitrarily weak
attraction only because there is a Fermi surface. The Higgs field must
condense in the vacuum.

\emph{Could the universe be filled with the medium?} No: a medium
defines a rest frame and would violate Lorentz invariance, and a Fermi
sea with the energy scale of the weak interaction would have an energy
density of order \((100\ \mathrm{GeV})^4\), some fifty-five orders of
magnitude above the density of the universe.

\emph{Can the pair condense without a medium?} Only if the attraction
exceeds a critical strength, as in the model of Nambu and Jona-Lasinio.
Without a Fermi surface the negative-energy states cannot be neglected,
and the gap equation of S1 does not apply. The appropriate equation is
the stationarity condition of the effective action of Cornwall, Jackiw
and Tomboulis for the fermion propagator in vacuum, with one exchanged
boson (the ``ladder'' approximation). S4, Sec. 7, derives and solves it
for the self-energy that a condensate of the vector pair channel
generates. The result (\texttt{data/cjt\_vacuum.csv}): a condensate
forms only if
\[\alpha_{VS}=\frac{g_{VS}^2}{4\pi}>\alpha_c\Big(\frac{\Lambda}{m},\frac{m_s}{m}\Big),\qquad
\alpha_c\to\frac{\pi}{2}\ \ (\Lambda\to\infty),\] where \(\Lambda\) is
the momentum at which the four-dimensional description ends:
\(\alpha_c=\) 3.74 for \(\Lambda=10\,m\), 2.15 for \(\Lambda=100\,m\)
and 1.64 for \(\Lambda=10^6\,m\) (at \(m_s=m\)). A vacuum condensate
therefore needs \(\alpha_{VS}\) of about 2 or more, that is
\(g_{VS}\gtrsim5\). At such couplings two fermions already form a bound
state (S1, Eq. (5.4)), the one-loop potential of the Vorton-Scalar is
unstable unless its self-coupling is equally strong (S4, Sec. 5.3), and
neither the one-boson-exchange kernel nor the ladder approximation is
quantitatively reliable. (An earlier draft of this section used the gap
equation of S1 with particles only and obtained a far larger critical
coupling; that estimate ignored the negative-energy states and is
superseded. Error log, E28.)

\textbf{Status of C2: open.} In the weak-coupling range, where the model
has been solved, there is no vacuum condensate. The ladder equation
indicates one for \(\alpha_{VS}\gtrsim2\), a strongly coupled regime in
which that equation is itself only a guide.

\begin{center}\rule{0.5\linewidth}{0.5pt}\end{center}

\hypertarget{conditions-c3-and-c4-the-electroweak-quantum-numbers}{%
\section{4. Conditions C3 and C4: the electroweak quantum
numbers}\label{conditions-c3-and-c4-the-electroweak-quantum-numbers}}

For a condensate of the field \(\psi\) to break the electroweak
symmetry, \(\psi\) must carry electroweak charges. The only symmetries
of the models of S1 that commute with the Lorentz group of ordinary
spacetime --- and that could therefore be internal symmetries of
four-dimensional physics --- are the rotations of the extra coordinates
and the two phase symmetries. The question is whether the electroweak
group fits into them, with the condensate in the right representation.

\hypertarget{the-71-model-the-pattern-of-the-higgs-vacuum}{%
\subsection{4.1 The 7+1 model: the pattern of the Higgs
vacuum}\label{the-71-model-the-pattern-of-the-higgs-vacuum}}

In 7+1 the four extra axes \(x0,x5,x6,x7\) are space-like, and their
rotation group \(SO(4)\), locally \(SU(2)_A\times SU(2)_B\), commutes
with ordinary rotations. For fields that do not depend on the extra
coordinates it is the complete non-Abelian internal symmetry (S1, Sec.
2.1; \texttt{data/reduced\_symmetry.csv}). The machine checks give the
representations:

\begin{itemize}
\tightlist
\item
  one particle at rest: \((2;2,1)\oplus(2;1,2)\) under
  \(SO(3)_{\rm ordinary}\times SU(2)_A\times SU(2)_B\) --- spin
  \(\tfrac12\), and a doublet of one of the two groups (check
  \texttt{71-Q1});
\item
  the s-wave pairs of ordinary spin 0: ten states,
  \((2,2)\oplus(3,1)\oplus(1,3)\) (check \texttt{71-Q7}); the \((2,2)\)
  are the four vector pairs \(\Delta^a\), \(a\in\{0,5,6,7\}\), a
  4-vector of \(SO(4)\);
\item
  the neutral current of the Vorton-Scalar attracts in the \((2,2)\) and
  repels in the \((3,1)\oplus(1,3)\) (S1, Sec. 4.2), and the state of
  lowest energy is \(\Delta^a=|\Delta|e^{i\varphi}n^a\) with a real unit
  vector \(n^a\) (S1, Sec. 6.3).
\end{itemize}

Compare with Sec. 1: \textbf{the order parameter is a \((2,2)\) of
\(SU(2)_A\times SU(2)_B\), and its vacuum value breaks
\(SO(4)\to SO(3)\) --- the representation and the breaking pattern of
the Higgs field and its custodial symmetry.} If \(SU(2)_A\) is
identified with \(SU(2)_L\) and the hypercharge with the generator
\(T^3_B\) of \(SU(2)_B\), the four components of \(\Delta^a\) have the
electric charges \(Q=T^3_A+T^3_B=(+1,0,0,-1)\), the hidden-axis
direction is electrically neutral, and conditions C3 and C4 are met.

Four qualifications:

\begin{enumerate}
\def\labelenumi{\arabic{enumi}.}
\tightlist
\item
  \emph{The order parameter is complex.} \(\Delta^a\) has four complex
  components --- twice the content of the Higgs doublet --- because the
  pair carries fermion number. Its vacuum value also breaks \(U(1)_F\).
  If fermion number is an exact global symmetry there is an additional
  massless Goldstone boson, and besides \(h\) a second neutral and a
  charged scalar; the structure is that of a model with two Higgs
  doublets (condition C7).
\item
  \emph{The identification of the gauge group is a hypothesis.} Nothing
  in S1 makes \(SU(2)_A\) a \emph{gauge} symmetry. In theories of the
  Kaluza--Klein type, rotations of extra dimensions become gauge
  symmetries when the extra space is curled up with the corresponding
  isometry group; a 4-torus (considered in the author's notebook
  \texttt{7+1.nb}, section 25) has only the isometries \(U(1)^4\), a
  4-sphere (its section 26) has \(SO(5)\supset SO(4)\). Whether Stage 2
  of the author's programme (the hidden axis curled up into a ring, the
  other three extra dimensions deflated) leaves an \(SU(2)\) is open; a
  ring by itself has only \(U(1)\), and treating \(x0\) differently from
  \(x5,x6,x7\) breaks the \(SO(4)\) that the custodial structure needs.
\item
  \emph{In the vacuum the multiplet is larger still.} A chemical
  potential singles out the pairs of fermion number 2. Without it, the
  zero-mode sector has further symmetries that mix particles with
  antiparticles (the full group is \(O(4)\times O(4)\); S4, Sec. 7.1),
  and the pair condensate is one orientation of a 16-component real
  multiplet that also contains particle--antiparticle condensates. A
  vacuum condensate would have to be analysed in that larger space.
\item
  \emph{Ordinary matter must carry the same charges.} Quarks and leptons
  are doublets of \(SU(2)_L\). If \(SU(2)_L\) is a rotation group of the
  extra dimensions, they must be spinors of those rotations,
  i.e.~components of fields like \(\psi\) itself. In the brane picture
  of \texttt{7+1.nb} (its sections 17 and 18) the Standard Model, with
  an elementary Higgs doublet, lives on a brane and has its own gauge
  groups; the field \(\psi\) couples to it through portals. That picture
  and the present hypothesis are alternatives, not parts of one model.
\end{enumerate}

\hypertarget{the-44-model-the-same-pattern-but-from-a-hidden-symmetry}{%
\subsection{4.2 The 4+4 model: the same pattern, but from a hidden
symmetry}\label{the-44-model-the-same-pattern-but-from-a-hidden-symmetry}}

In 4+4 the hidden axis is space-like and the three other extra axes are
time-like. Two symmetry groups must be distinguished (S1, Sec. 2.1).

\emph{The symmetry that the geometry provides.} The rotations of the
three additional times, \(SO(3)_T\), are a symmetry of the quantized 4+4
model whether or not the fields depend on \(x0\). Under \(SU(2)_T\) one
particle at rest is a doublet (check \texttt{44-Q1}); pairs are singlets
or triplets (check \texttt{44-Q2}); the hidden-axis pair \(\Delta^0\) is
a singlet and \(\Delta^5,\Delta^6,\Delta^7\) form a triplet (check
\texttt{44-Q3}). If the weak isospin were \(SU(2)_T\), the hidden-axis
condensate would not break the electroweak symmetry at all, and a
condensate of the triplet would break it with \(\rho\ne1\) at tree level
(\(\rho=\tfrac12\) for hypercharge 1; no \(Z\) mass at all for
hypercharge 0), against the measured value 1 to a part in a thousand.

\emph{The hidden symmetry of the zero-mode sector.} For fields that do
not depend on \(x0\) either, the model has the larger symmetry
\(SU(2)_A\times SU(2)_B\), of which \(SU(2)_T\) is the diagonal
subgroup. The particles are \((2,1)\oplus(1,2)\), the four pairs
\((\Delta^0,i\Delta^5,i\Delta^6,i\Delta^7)\) form the \((2,2)\), and the
condensate breaks \(SO(4)\to SO(3)_T\): algebraically the same pattern
as in 7+1. But the three generators beyond \(SO(3)_T\) are the matrices
\(\gamma^0\gamma^t\), which act on the components of the field only.
They are not rotations of the spacetime --- the transformations of the
spacetime that mix \(x0\) with a time are boosts, generated by the same
matrices \emph{without} the factor \(i\), and boosts are not unitary.
Hence:

\begin{itemize}
\tightlist
\item
  the hidden symmetry is not an isometry of any extra space, so the
  mechanism of Kaluza and Klein cannot make it a gauge symmetry;
\item
  it is violated by every mode of the fields that depends on \(x0\);
\item
  it could be gauged only by hand, in the four-dimensional theory of the
  zero modes --- a new postulate, unrelated to the geometry of 4+4.
\end{itemize}

\textbf{In 4+4 the algebraic pattern of the Higgs vacuum is present in
the zero-mode sector, but the electroweak gauge group has no geometric
origin there.} This is the precise form of the difference between the
two spacetimes: the metric on the four extra axes is
\(\mathrm{diag}(-,-,-,-)\) in 7+1 and \(\mathrm{diag}(+,-,-,-)\) in 4+4,
the compact group \(SO(4)\) is the isometry group only of the first, and
in the second it survives as a symmetry of the zero modes alone.

\begin{center}\rule{0.5\linewidth}{0.5pt}\end{center}

\hypertarget{conditions-c5-c6-and-c7-scale-masses-spectrum}{%
\section{5. Conditions C5, C6 and C7: scale, masses,
spectrum}\label{conditions-c5-c6-and-c7-scale-masses-spectrum}}

These conditions concern dynamics that the present model does not
determine. What can be said:

\textbf{C5, scale and mass of the scalar particle (estimate).} In models
of the Nambu--Jona-Lasinio type the scalar bound state that accompanies
the condensate has a mass of about twice the mass that the condensate
gives to its constituent fermions. With 125 GeV for the scalar, the
constituents would weigh about 60 GeV. By C3 they carry electroweak
charges; particles of that mass with electroweak charges would have been
produced in pairs at the electron--positron collider LEP, whose searches
exclude new charged fermions below about 100 GeV, and in decays of the
\(Z\) boson if lighter than 45 GeV. Refinements of the relation between
the two masses (Bardeen, Hill and Lindner) change it logarithmically,
not by the factor needed. \emph{This estimate speaks against the
hypothesis}; it is not a proof, because the strong- coupling dynamics of
Sec. 3 is not under control.

\textbf{C6, masses of quarks and leptons.} The model contains no
coupling of the pair field to ordinary fermions. A fermion-number-2 pair
cannot couple to a fermion--antifermion bilinear of ordinary matter
without violating either fermion number or the conservation of the
ordinary quantum numbers; an additional mechanism would have to be
invented. \emph{Not addressed.}

\textbf{C7, additional particles.} By Sec. 4.1 the 7+1 model predicts,
besides the scalar \(h\): the partners that complete the complex
\((2,2)\) --- a second neutral and a charged scalar --- a Goldstone
boson of fermion number, the fermions \(\psi\) themselves with
electroweak charges, and the two modes of the Vorton-Scalar. None of
these is observed. \emph{An obstacle, of a size that depends on unknown
masses.}

\textbf{The three hypotheses of the task need three different scales.}
The Higgs hypothesis needs the condensate at 246 GeV. The dark-matter
analysis of S3 points to bound pairs heavier than 0.2 GeV with a
nuclear-scale coupling, or to no superfluid at all. The Hubble-tension
mechanism of S3 needs fermions of about 1 eV. One field with one mass
cannot serve all three purposes.

\begin{center}\rule{0.5\linewidth}{0.5pt}\end{center}

\hypertarget{verdict}{%
\section{6. Verdict}\label{verdict}}

\begin{longtable}[]{@{}
  >{\raggedright\arraybackslash}p{(\columnwidth - 6\tabcolsep) * \real{0.2500}}
  >{\raggedright\arraybackslash}p{(\columnwidth - 6\tabcolsep) * \real{0.2609}}
  >{\raggedright\arraybackslash}p{(\columnwidth - 6\tabcolsep) * \real{0.2609}}
  >{\raggedright\arraybackslash}p{(\columnwidth - 6\tabcolsep) * \real{0.2283}}@{}}
\toprule\noalign{}
\begin{minipage}[b]{\linewidth}\raggedright
Condition
\end{minipage} & \begin{minipage}[b]{\linewidth}\raggedright
7+1
\end{minipage} & \begin{minipage}[b]{\linewidth}\raggedright
4+4 (quantized sector)
\end{minipage} & \begin{minipage}[b]{\linewidth}\raggedright
Basis
\end{minipage} \\
\midrule\noalign{}
\endhead
\bottomrule\noalign{}
\endlastfoot
C1 scalar in ordinary spacetime & met & met & DM Sec. 3.3; checks
\texttt{71-K3}, \texttt{44-K3} \\
C2 condensate in the vacuum & open (needs \(\alpha_{VS}\gtrsim2\)) &
open (needs \(\alpha_{VS}\gtrsim2\)) & ladder equation, S4 Sec. 7 \\
C3 electroweak doublet with neutral component & met, \emph{if}
\(SU(2)_L=SU(2)_A\) and \(Y=T^3_B\) & only under a hidden symmetry that
the geometry cannot gauge; singlet under \(SO(3)_T\) & checks
\texttt{71-Q1}, \texttt{71-Q7}, \texttt{44-Q1} to \texttt{44-Q3}; S1
Sec. 2.1 \\
C4 breaking \(SO(4)\to SO(3)\) & met & met in the zero-mode sector only
& S1 Secs. 2.1, 6.3 \\
C5 scale and scalar mass & unfavourable (estimate) & --- & Sec. 5 \\
C6 masses of ordinary fermions & not addressed & --- & Sec. 5 \\
C7 no unobserved particles & additional particles predicted & --- & Sec.
5 \\
\end{longtable}

\textbf{What has been established.} (i) In both spacetimes the
hidden-axis pair is a scalar of ordinary spacetime. (ii) In 7+1 it
belongs to a \((2,2)\) of the rotation group of the extra space axes,
and its condensate has the custodial breaking pattern of the Higgs
vacuum. (iii) In 4+4 the same pattern exists for the zero modes, as a
hidden symmetry; under the rotations of the additional times, the only
non-Abelian internal symmetry that the geometry of 4+4 provides, the
condensate is a singlet.

\textbf{What has not been established.} That a vacuum condensate forms;
that rotations of the extra dimensions are the electroweak gauge group;
that the scale, the mass of the scalar and its couplings come out as
observed; that ordinary fermions obtain their masses. The hypothesis
``the singlet superfluid is the Higgs field'' is therefore
\textbf{unproven in 7+1 --- with one structural point in its favour and
two estimates against it --- and without a geometric basis in 4+4},
where the structural point exists only for a hidden symmetry.

A remark on what the model \emph{does} contain: the Vorton-Scalar
\(\Phi_{VS}\) is itself a scalar field with a vacuum value that gives
the fermions \(\psi\) their mass (S1, Eq. (1.5)). It plays for the dark
sector the role that the Higgs field plays for ordinary matter, but it
is a singlet of all non-Abelian symmetries and cannot be the Higgs field
of the Standard Model either.

\hypertarget{references}{%
\section{References}\label{references}}

\begin{enumerate}
\def\labelenumi{\arabic{enumi}.}
\tightlist
\item
  P. W. Higgs, \emph{Broken symmetries and the masses of gauge bosons},
  Phys. Rev.~Lett. \textbf{13}, 508 (1964); F. Englert and R. Brout,
  \emph{Broken symmetry and the mass of gauge vector mesons}, Phys.
  Rev.~Lett. \textbf{13}, 321 (1964).
\item
  S. Weinberg, \emph{A model of leptons}, Phys. Rev.~Lett. \textbf{19},
  1264 (1967).
\item
  P. Sikivie, L. Susskind, M. Voloshin and V. Zakharov, \emph{Isospin
  breaking in technicolor models}, Nucl. Phys. B \textbf{173}, 189
  (1980).
\item
  Y. Nambu and G. Jona-Lasinio, \emph{Dynamical model of elementary
  particles based on an analogy with superconductivity. I}, Phys.
  Rev.~\textbf{122}, 345 (1961).
\item
  W. A. Bardeen, C. T. Hill and M. Lindner, \emph{Minimal dynamical
  symmetry breaking of the standard model}, Phys. Rev.~D \textbf{41},
  1647 (1990).
\item
  S. Navas et al.~(Particle Data Group), \emph{Review of particle
  physics}, Phys. Rev.~D \textbf{110}, 030001 (2024), and 2025 update
  (values of \(G_F\), \(m_H\), \(\rho\); limits on heavy charged
  leptons).
\item
  E. Witten, \emph{Search for a realistic Kaluza--Klein theory}, Nucl.
  Phys. B \textbf{186}, 412 (1981).
\item
  J. M. Cornwall, R. Jackiw and E. Tomboulis, \emph{Effective action for
  composite operators}, Phys. Rev.~D \textbf{10}, 2428 (1974).
\item
  P. L. Nash, notebook \texttt{7+1.nb}, sections 17, 18, 25, 26 (see
  \texttt{provenance/P3\_notebooks\_used}).
\end{enumerate}

\end{document}
